\section{Methods}

\subsection{Governing equations and nondimensionalization}

We considered an infinite homogeneous elastic plate occupying
\(\mathbb R^2\times[-h,h]\), with density \(\rho\), elasticity tensor
\(C_{ijkl}\), and traction-free faces.  For a plane-wave component
\(\boldsymbol u=\boldsymbol U(z)\exp[\mathrm i(k_xx+k_yy)-\mathrm
i\omega t]\), the strong problem is
\(\partial_j(C_{ijkl}\partial_lu_k)=\rho\,\partial_t^2u_i\), with
\(\boldsymbol\sigma\boldsymbol n=0\) at \(z=\pm h\).  The reference
isotropic solid has \(\lambda=2\), \(\mu=1\), \(\rho=1\), and \(h=1\).
The shear speed is \(c_T=\sqrt{\mu/\rho}\); reported wavenumber, frequency,
and time are \(\kappa=h k^{\rm phys}\),
\(\Omega=h\omega^{\rm phys}/c_T\), and \(t=c_Tt^{\rm phys}/h\).
Consequently, every plotted coordinate and fitted time is dimensionless.
The signed parameter \(\varepsilon\) multiplies the declared cubic
constitutive perturbation at fixed density; it is not a frequency shift.

The isotropic symmetric branch was first computed from the desingularized
Rayleigh--Lamb determinant, rather than from a plate approximation.  Its first
\(S_1\)-ZGV root in the declared search window was obtained at 80-decimal
working precision from the
simultaneous determinant and radial-derivative equations.  Implicit
differentiation supplied the radial curvature.  The selected root and
curvature then served as independent regression targets for the
thickness-resolved solver.  The full determinant construction,
desingularization at the shear cutoff, and Morse--Bott Hessian calculation are
given in Supplementary Section~\ref{sec:isotropic-rayleigh-lamb-derivation}.

\subsection{Spectral solvers}

The three displacement components were discretized only through the
thickness by a Legendre--Gauss--Lobatto (GLL) spectral element method.  The
weak form retains the face traction as its natural boundary term; no face
degree of freedom is constrained and no penalty is added.  Stiffness and mass
were assembled by collocated GLL quadrature in the Mandel strain convention.
The production homogeneous-plate mesh uses one element on \([-h,h]\); a
two-element mesh split at \(z=0\) is an independent assembly check.  The
generalized Hermitian problem
\(\boldsymbol K\boldsymbol u=\omega^2\boldsymbol M\boldsymbol u\) was solved
for the lowest 12 modes with SciPy's generalized-Hermitian \texttt{gvx}
driver.  Eigenvectors were mass normalized.  Matrix Hermiticity, mass
orthogonality, and the normalized eigen residual were tested at every declared
convergence order.

All first and second derivatives of \(\boldsymbol K\) with respect to
\((k_x,k_y)\) were assembled analytically from the wavevector-affine strain
matrix.  Absolute frequency and gradient discrepancy estimators were the
differences between answers at orders \(p\) and \(p+4\).  These inter-order
quantities are empirical resolution indicators, not bounds on the unknown
continuum errors.  The relative estimator used by the eigengap gate was the
maximum of that nested-order eigenvalue discrepancy and the two eigen
residuals.  The committed full profile uses order
pairs 12/16 for the isotropic artifact, 10/14 for the critical-point, scaling,
and Green artifacts, and 12/16 for the silicon stress test; convergence sweeps
the pairs 6/10 through 14/18.  The sensitivity artifact uses order 10 and is
checked by independent physical finite differences rather than a 10/14 pair.  The exact
weak form, derivative matrices, residual definitions, and all fixed gates are
specified in Supplementary Section~\ref{sec:spectral-numerics}.

\subsection{Mode tracking and sensitivities}

Parity identifies the symmetric branch in the independent isotropic validation.
Production angular anchors were seeded by proximity to the exact target
frequency and then continued by a mass-weighted modal assurance criterion.  Near a close pair,
tracking used the complete connected eigenvalue cluster and its principal
subspace overlap before phase alignment.  Eight endpoint-free angular anchors
were constructed independently at orders \(p\) and \(p+4\); closure of the
angular loop required both scalar MAC and the least principal cosine to be at
least \(1-10^{-8}\).  Inward and outward radial rays were continued separately
from each ring anchor.  Samples were rejected when the selected cluster
touched the top of the 12-mode solve or when the relative eigengap was no more
than ten times the relative eigenvalue discrepancy estimator.  This eigengap
gate is a method hypothesis: no scalar sensitivity or Hessian was accepted
through a numerically unresolved crossing.

The generalized-eigenvalue sensitivity was evaluated at fixed mass norm by
the Hellmann--Feynman contraction.  The angular function \(V(\theta)\) and its
mean and fourth harmonic, \(V_0\) and \(V_4\), were computed on 64
endpoint-free angles.  The mixed radial coefficient \(B\) retained the
differentiated-mode contribution.  The mode derivative was obtained from a
bordered reduced-resolvent system with the mass-gauge constraint, and both the
equation and gauge residuals were required to be below \(10^{-8}\).  A centered
full-wave difference in \(\varepsilon\) checked \(V(\theta)\), and a mixed
\((\varepsilon,q)\) difference checked \(B(\theta)\), both with
\(\Delta\varepsilon=\Delta q=0.0025\).  An independent tensor-product step sweep
\(0.02,0.01,0.005,0.0025\) checked convergence.  Supplementary
Section~\ref{sec:anisotropic-morse-derivation} gives the complete
generalized-pencil algebra, rotated cubic contraction, normal form, and Morse
splitting proof.

\subsection{Critical-point search and annular consistency check}

Critical points were searched only on the tracked branch inside the annulus
\((1-f_A)k_0\leq|\boldsymbol k|\leq(1+f_A)k_0\), with \(f_A=0.15\).
Every polar cell in which both radial and tangential gradient components
bracketed zero supplied a seed.  Each seed was refined in Cartesian
coordinates by \texttt{scipy.optimize.root} with the default Powell hybrid
method, a Richardson Cartesian Hessian from steps \(h_H\) and \(h_H/2\) as
Jacobian, and solver tolerance
\(10^{-11}\).  A root outside the annulus or with a gradient residual exceeding
its nested-order gradient-discrepancy estimator was rejected.  Roots were
merged only within the propagated, capped position-resolution radius.

The full calculation repeated the candidate search on \(5\times16\) and
\(9\times32\) radial-by-angular grids.  The Hessian was Richardson extrapolated
from centered gradient differences with full-profile step \(10^{-3}\); a
classification was accepted only if its least absolute eigenvalue exceeded
ten times the Hessian discrepancy estimate.  Counts, kinds, locations,
frequencies, and Hessian eigenvalues had to agree under candidate-grid
doubling.  An
independent annular check sampled each boundary at 32 and 64 angles, required a
gradient norm separated from its discrepancy estimate, stable boundary
winding, angular increments below \(\pi/2\), and equality between outer-minus-inner
winding and the sum of point indices.  The \(\varepsilon<0\) search was repeated
on an angularly offset candidate grid.  This is a finite-resolution
local-annulus consistency check, not an exhaustive continuum enumeration and
not a global enumeration of every Lamb and shear-horizontal branch.  Index
closure alone also cannot exclude an unresolved pair of
opposite-index critical points between candidate-grid nodes; candidate-grid
doubling, symmetry-offset repetition, and Hessian separation supply the
additional controls used here.

A separate finite-anisotropy stress test used crystal axes
\(x\parallel[100]\), \(y\parallel[010]\), and \(z\parallel[001]\), so the
plate normal was [001].  The constants \(C_{11}=165.7\), \(C_{12}=63.9\), and
\(C_{44}=79.6\) GPa followed the published single-crystal elasticity record
\citep{zhang2014silicon}.  The calculation normalized stiffness by \(C_{44}\),
set \(\rho=h=1\), and reported frequency in
\(\sqrt{C_{44}/\rho}/h\).  It exercised the search and annular consistency check
under stronger constitutive contrast and was not used to validate the
weak-\(\varepsilon\) theorem.

\subsection{Green-function quadrature and fit masks}

The positive-frequency response was projected onto the tracked branch for a
normal impulsive source and top-face normal detector.  The modal amplitude
was \(\mathrm i|u_z(h)|^2/(2\omega)\).  A Gaussian source factor with
\(R/h=0.5\) and a fixed branch window
\(\exp[-(q/\sigma)^8]b_\sigma(q)\), where \(q=|\boldsymbol k|-k_0\) and
\(\sigma/k_0=0.15\), were applied.  The deterministic \(C^\infty\) taper
\(b_\sigma\) equals one for \(|q|\leq1.25\sigma\), uses the standard flat
\(\exp(-1/x)\) transition, and is exactly zero for
\(|q|\geq1.5\sigma\).  Thus the integration interval
\([-1.5\sigma,1.5\sigma]\) contains the complete compact support and has no
endpoint contribution.  Direct integration used the Cartesian
Fourier measure with radial polar Jacobian \(k_0+q\), an endpoint-free
trapezoidal rule in angle, and composite
Simpson quadrature in radius.  No Bessel function entered this full-wave
quadrature.  The full profile rebuilt independent tracked surfaces on
\(129\times32\), \(257\times64\), and \(513\times128\) grids for the uniform
comparisons at \(\varepsilon\leq0.04\); the fixed \(\varepsilon=0.08\)
comparison used twice as many angular nodes.
For every response row, the stationary-point search used an annular half-width
\(1.55\sigma=0.2325k_0\), strictly larger than the direct support
\(1.5\sigma=0.225k_0\), on a \(15\times32\) candidate grid.  Boundary/index
checks used 16 and 32 angular nodes.  Every resolved stationary point in this
larger domain was evaluated with the same compact taper and included in the
fixed-Morse sum; the domain, margin, point counts, and contribution counts are
stored in the Green sidecar.
The computed time axis was \(t_n=400+25n\), \(0\leq n\leq504\), for
\(\varepsilon\in\{0.005,0.01,0.02,0.04,0.08\}\).

Fit masks were fixed in code before response evaluation.  The early slope
used \(\varepsilon=0.005\), \(t\geq1500\), and
\(0.10\leq|\varepsilon V_4t|\leq0.30\).  The fixed-anisotropy late slope used
complete half-open minimum--saddle modulation periods beginning at \(t=1500\),
with at least three samples per period and four periods.  Ordinary least
squares was applied to \(\log t\) and the logarithm of the predeclared RMS or
demodulated amplitude.  The fitted intercepts were diagnostic only: the Bessel and
exact-Morse comparisons used no fitted alignment of amplitude, phase, carrier
frequency, or time.  The full response was accepted only if the three-grid
complex error contracted and the accumulated inter-resolution
phase-discrepancy estimator satisfied the declared
\(t_{\max}\max|\delta\omega|\leq0.05\) rad acceptance test.  This estimator
compares computed resolutions and does not bound the unknown continuum phase
error.  Supplementary
Section~\ref{sec:green-function-asymptotics} derives the normalized response
and all asymptotic regimes; Supplementary Section~\ref{sec:spectral-numerics}
specifies the quadrature, fit masks, cancellation exclusion, and error gates.

\subsection{Reproducibility}

Reference physical parameters and cross-stage acceptance limits are versioned
in \path{config/reference.yaml}; stage-specific grids, masks, and thresholds
are versioned in the workflow source and repeated in artifact sidecars.  The
interpreter is pinned by \path{.python-version}, and dependency resolution is
fixed by \path{uv.lock}.  The seven workflow stages write paired NPZ arrays and
JSON sidecars containing array schemas, units, commands, input hashes, source
and lock hashes, tolerances, measured diagnostics, and SHA-256 output hashes.
Figures read only qualified artifacts, export their plotted source CSV files,
and pass a strict geometry, provenance, and typography audit.

The committed environment used uv 0.11.19, Python 3.12.13, NumPy 2.5.1, SciPy
1.18.0, SymPy 1.14.0, Matplotlib 3.11.0, mpmath 1.3.0, and PyYAML 6.0.3.
Supported Python versions and dependency ranges remain declared in
\path{pyproject.toml}; the interpreter pin and lock file, not those ranges,
define the verified Python environment.  PDF compilation used latexmk 4.88,
pdfTeX 1.40.29, and BibTeX 0.99e from TeX Live 2026; these are external build
prerequisites rather than Python dependencies.

From the package root, the currently executable scientific rebuild is

\begin{verbatim}
uv sync --python 3.12.13 --frozen --all-extras
uv run python scripts/reproduce_all.py \
  --profile full --skip-paper
uv run python scripts/export_manuscript_values.py
\end{verbatim}

The full-profile command first runs the complete test suite, then the
isotropic, sensitivity, critical-point, scaling, Green, convergence, and
silicon stages in isolated processes.  It validates each artifact immediately,
validates the provenance manifest, rebuilds the six main and six supplementary
figures and two generated tables, and performs strict figure QA.  The next
command exports the evidence-checked manuscript macros.  The locked compiler
then checks the recorded latexmk, pdfTeX, and BibTeX versions, builds the
Supplement from \path{supplement.tex} before the main text from
\path{main.tex} in two
independently cleaned cycles, applies the strict final-log gate, rejects
nondeterministic PDF metadata, and requires the two cycles to agree byte for
byte:

\begin{verbatim}
uv run python scripts/compile_paper.py
\end{verbatim}

The supplement must be compiled first because the main document imports its
label file.  The smoke profile is a runtime check only; it is not the source of
committed numerical evidence.  Availability boundaries are stated below.
